\paragraph{The cyclotomic level models.}\label{par:wach-level-models}
Fix the same crystalline lattice map \(f:T_1\to T_2\) as in
Section~\ref{sec:inert-long-orbit}. Write \(K_O=\operatorname{Frac}(O)\)
for the coefficient field, and retain the choice containing the
cyclotomic roots. Consider an arc of length \(a\) from \(c-a\) to
\(c\). At backward offset \(s\), the letter is \(c-s\). A jump of
the decreasing Hodge filtration at \(u\) means that
\(\operatorname{Fil}^u\ne\operatorname{Fil}^{u+1}\).
Let \(G_0\le g\) bound the absolute jumps at thin letters. We
can take \(G_0=g\) for the given thin box.

Write \(b_s=s_{c-s}\)
for the size of a position with larger jumps. Upper and lower bounds \(m_s^+\)
and \(-m_s^-\) for the jumps, and the exponent of the arc
denominator, are
\[
\begin{array}{c|ccc}
 & m_s^+ & m_s^- & H_s^{\mathrm{arc}}=g+m_s^+\\ \hline
 \text{thin} &G_0&G_0&g+G_0\\
 \text{up}   &b_s+S_0&\eta_0&g+b_s+S_0\\
 \text{down} &\eta_0&b_s+S_0&g+\eta_0 .
\end{array}
\]
Here up means that the large jumps are in \([b_s,b_s+S_0]\).
For down they are in \([-b_s-S_0,-b_s]\).
The arc exponents are used in
\[
 \Lambda=\prod_{s<a}p_s^{H_s^{\mathrm{arc}}},\qquad
 \Theta=\prod_{s<a}p_s^g .
\]
The symmetric widths \(H_i\) from the recurrence for \(R\) control
the ordinary Wach cokernel.

For the evaluation at offset \(s\), use the product over all
primitive \(\ell^{s+1}\)-th roots \(\zeta\) of the space of Laurent
series in \(t\) with coefficients in the terminal crystalline
space \(D_{i,c}\). In each component expand the ordinary
cyclotomic parameter at \(1+\pi=\zeta\exp(t)\).
The scalar evaluation map is denoted by \(\rho_s\).

At every primitive root indexing \(\rho_s\) there is an expansion
\[
 \rho_s(p_s)=t\,u_s(t),\qquad u_s(0)\ne0.
\]
Indeed the primitive roots are simple
roots of the cyclotomic polynomial in characteristic zero and
\(\zeta\exp(t)\) has a nonzero derivative at zero.
The tame parameter \(X\) and all \(p_{s'}\) with \(s'\ne s\)
evaluate to units. Frobenius on the series sends
\(\zeta\) to \(\zeta^\ell\) and rescales \(t\) by \(\ell\),
and on the crystalline coefficients it transports them to the
next letter. Cyclotomic action on the degree \(k\) coefficients
is the permutation of the roots given by the cyclotomic value
times its \(k\)-th power. The level space consists of the tuples
fixed by the tame group for this action.

If \(\operatorname{Fil}^u_{c-s}\) denotes the filtration at the starting letter, the filtration at offset \(s\) is
\[
 \mathcal F_{i,s}^u=\Phi_{\rm cris}^s(\operatorname{Fil}^u_{c-s})\subseteq D_{i,c}.
\]
This transport is the de Rham comparison at that level. We label a Wach
step by its terminal Hodge letter. To translate
the de Rham convention using \(\varphi^{-(s+1)}\), label a
Wach component by the embedding one step before its
ordinary coefficient embedding and transport its crystalline
coefficients by \(\Phi_{\rm cris}^{-1}\). An evaluation at
level \(s=0\) then has the Hodge filtration of its own letter.
At level \(s\) the remaining \(s\) steps give the filtration of
\(c-s\), transported by \(\Phi_{\rm cris}^s\).
This convention is common to the source, target and extension.
The Wach evaluation generates the completed rational module
of tuples having coefficient of degree \(k\) in
\(\mathcal F_{i,s}^{-k}\). We denote this condition by
\(J_{i,s}\).

\begin{lemma}[The completed Wach lattice at a cyclotomic root]\label{lem:wach-evaluated-hodge}
Label a Wach step by its terminal Hodge letter and use
\(\Phi_{\rm cris}^{s}\) to transport the filtration from \(c-s\).
If \(\widehat N_{i,s}^{\rm rat}\) denotes the completed rational Wach lattice at cyclotomic level \(s\), then
\[
 \widehat N_{i,s}^{\rm rat}=J_{i,s}\qquad(0\le s<a),
\]
with \(J_{i,s}\) as in Paragraph~\ref{par:wach-level-models}. At the starting letter
reduction at \(X=0\), after inverting \(\ell\), gives the crystalline
space. Both comparisons commute with the coefficient-field action.
\end{lemma}

\begin{proof}
First use a common Tate twist to put the
crystalline weights in the positive range of Berger's Wach
comparison \cite[Propositions II.2.1, II.3.1 and III.2.1]{berger-wach}.
Over the ordinary cyclotomic parameter, the comparison at a
root is the de Rham embedding indexed by that root, with
\(t\) rescaled and crystalline coefficients transported by
Frobenius. If \(d_j\) is an adapted basis with jumps \(u_j\), the completed de Rham lattice is
\[
 \bigoplus_j K_O[[t]]\,t^{-u_j}d_j
\]
in a root component with coefficient field \(K_O\). Indeed these elements lie in the de Rham
lattice and their reductions form a basis by the Hodge--Tate
comparison for the de Rham representation. The completed
Wach module maps into this lattice through
\(N(V)\subset D^+(V)\) in the positive range. Relative to the
lattice generated by the constant crystalline basis, its
elementary-divisor exponents at the root are \(u_j\) by
Berger's two elementary-divisor propositions.

Let \(L\subseteq L'\) denote the inclusion at the root. Their determinant orders relative to the same constant basis are both \(\sum_j u_j\), so
\[
 \operatorname{length}(L'/L)
  =\operatorname{ord}\det L-\operatorname{ord}\det L'=0.
\] This comparison
can first be made on the underlying \(\Q_\ell\)-representation.
It commutes with the coefficient-field action and hence with
the idempotents splitting the coefficients and the embeddings.

Undoing the Tate twist shifts the jumps and the Laurent orders
by the same amount, since the change of Wach basis and of
crystalline basis differ by a power of \(t/\pi\). At a primitive root \(\zeta\ne1\) the local parameter \(t\) vanishes and \(\pi(0)=\zeta-1\ne0\). Hence
\[
 \operatorname{ord}_t(t/\pi)=1.
\]
This is the shift in Laurent order when the Tate twist is undone.

Passage to tame invariants
uses the averaging idempotent for the group of order
\(\ell-1\). Averaging the ordinary basis preserves its reduction
modulo \(\pi\), where the action is trivial. The averaged basis
therefore remains a basis over \(O[[\pi]]\), and its invariant
coefficients are in \(\mathcal R\). This yields the asserted
coefficient condition on the tame-invariant spaces. At the root one, the starting-letter identification is
\[
 N_{i,c-a}[1/\ell]/XN_{i,c-a}[1/\ell]\simeq D_{i,c-a}.
\]
It is the crystalline comparison at the origin
\cite[Corollary III.4.5]{berger-wach}. In crystalline
coordinates the ordinary expansion at one is regular, and
the factor \(t/\pi\) for the twist is a unit there.
\end{proof}

\begin{proposition}[An integral membership test along an arc]\label{prop:wach-arc-membership}
Write \(w_s\) for the level-\(s\) evaluation and put
\[
 W_i=\Lambda^{-1}\operatorname{span}_{\mathcal R}\Phi^a N_{i,c-a}.
\]
Using \(m_s^\pm\) and \(\Lambda\) from
Paragraph~\ref{par:wach-level-models}, one has
\[
 \Bigl(\prod_{s<a}p_s^{m_s^+}\Bigr)N_{i,c}
    \subset\operatorname{span}_{\mathcal R}\Phi^aN_{i,c-a},
 \qquad
 w\in N_{i,c}\ \Longleftrightarrow\
 (\forall s<a)\ w_s\in J_{i,s}\quad(w\in W_i).
\]
The signed step bound applies also to the crystalline extension of
\(O\) by \(T_2\).
\end{proposition}

\begin{proof}
\proofstep{Clearing the one-step denominators.}
These comparisons give integral step and level tests. For
one step ending at \(i\), the span \(\mathcal R\Phi N_{i-1}\)
evaluates to the regular lattice at the first cyclotomic
level. This follows by transport of the comparison at one.
The condition at that level on \(N_i\) is the Hodge filtration
at the ending letter. At the first level the two factors move in opposite directions:
\[
 q^{m^+}N_i\ \subseteq\ \mathcal R\Phi N_{i-1},\qquad q^{m^-}\mathcal R\Phi N_{i-1}\ \subseteq\ N_i,
\]
initially in the evaluated modules. For integral inclusions use the coarse Wach height bounds.
After a common twist the ordinary positive Wach module is
stable under Frobenius, with linearized cokernel killed by
the cyclotomic factor to the power of the width of the
jump box \cite[Definition III.4.1 and Proposition III.4.2]{berger-wach}.
Undoing the twist multiplies the Frobenius matrix by a
power of that factor since \(\phi(\pi)/\pi\) is the ordinary
factor.

Tame averaging gives the corresponding inclusions
on invariant modules. Thus the coordinates on either basis
have an a priori denominator a power of \(q\). The
local test removes that denominator integrally. In fact the
remainder on division by the distinguished associate of
\(q^m\) vanishes when its Taylor coefficients of order
less than \(m\) vanish at the roots. Indeed let \(D\) be the distinguished polynomial and \(R\) the remainder. Over \(K_O\) the Taylor conditions imply
\[
 D\mid R,\qquad \deg R<\deg D\quad\Longrightarrow\quad R=0.
\]
The quotient by \(D\) is free over \(O\), so this vanishing descends integrally. Iterating the resulting signed bound gives
\[
 \Bigl(\prod_{s<a}p_s^{m_s^+}\Bigr)N_{i,c}
    \subset\operatorname{span}_{\mathcal R}\Phi^aN_{i,c-a}.
\]
The same proof applies to the crystalline extension
\(0\to T_2\to T'\to O\to0\). Exactness of the crystalline filtration
adds only the jump at zero, which lies in each of the stated boxes.

\proofstep{Testing the product of the denominators.}
Put
\[
 W_i=\Lambda^{-1}\operatorname{span}_{\mathcal R} \Phi^aN_{i,c-a}.
\]
Using the opposite step bound, there is a product
\(Q_*\) of powers of the factors \(p_s\) with \(Q_*W_i\subset N_{i,c}\).
For \(w\in W_i\), expand \(Q_*w\) on a Wach basis at \(c\).
At each level the Wach basis is a basis of the completed Hodge
lattice at each root. Thus the evaluation of \(w\) is in \(J_{i,s}\)
exactly when its multiplied coordinates for \(Q_*w\) vanish to the
order of \(\rho_s(Q_*)\) at every root of level \(s\).

For the scalar
remainder test over \(O[[X]]\), the values of \(X\) at primitive
roots identify the distinct tame orbits, with one simple root of the
distinguished associate of \(p_s\) per orbit. Weierstrass division
and the same root test therefore give divisibility by the entire
product \(Q_*\) in every coordinate. Consequently
\[
       w\in N_{i,c}\quad\Longleftrightarrow\quad
       w_s\in J_{i,s}\text{ for every }s<a .
\]
The units at each level give the two tests used below. The inverse
image of \(N_{i,c}\) under \(\Theta\) lies in \(W_i\) by the signed
step bound. In \(W_i\) its test is \(p_s^g w_s\in J_{i,s}\)
for every \(s\). Also \(XN_{i,c}\) consists of the vectors of
\(XW_i\) having evaluations in every \(J_{i,s}\); apply the
previous test after dividing by \(X\), a unit at the evaluated
levels.
\end{proof}

Define \(D_{i,s}\) to be the diagonal crystalline constants.
For an up letter write \(U^K_{i,s}=\mathcal F_{i,s}^{b_s}\),
and for down write \(U^K_{i,s}=\mathcal F_{i,s}^{-\eta_0}\).
Define \(B_{i,s}\) by the following increasing degree filtration
on its coefficients. In a thin letter it is zero below degree
zero and full from zero. In an up letter the breaks are \(-b_s\)
and zero. In a down letter they are zero and \(b_s\). In the
last two cases it is zero before the first break, \(U^K_{i,s}\)
from the first to the second break, and full from the second.
Each coefficient condition is imposed at all the roots.

\begin{lemma}[Comparison with the rounded degree filtration]\label{lem:wach-rounded-breaks}
For the coefficient spaces \(B_{i,s}\) with the rounded breaks
just defined, one has
\[
 p_s^gB_{i,s}\subset J_{i,s},\qquad p_s^gJ_{i,s}\subset B_{i,s}.
\]
In coordinates the rounded filtration gives the divisibility conditions
of \eqref{eq:wach-level-condition}, compatibly with the crystalline map.
\end{lemma}

\begin{proof}
Between the two jump windows the Hodge filtration is constant,
with middle subspace precisely \(U^K_{i,s}\). Its two groups
of degrees differ from the rounded breaks by at most \(g\)
since \(G_0,S_0,\eta_0\le g\). In the thin case the one
group differs from zero by at most \(G_0\). It follows,
using multiplication by the unit part of \(\rho_s(p_s)\) on
the increasing filtrations, that
\[
       p_s^g B_{i,s}\subset J_{i,s},\qquad
       p_s^g J_{i,s}\subset B_{i,s}.
\]

Fix a Wach basis \(e_i\) at the start and write \(\lambda_i\)
for its reduction at one transported by \(\Phi^a\). This is
a basis of \(D_{i,c}\). Viewing the basis as a row, the coordinate evaluation at offset \(s\) is
\[
 y\longmapsto\rho_s(\Lambda)^{-1}\lambda_i\rho_s(y),\qquad y\in\mathcal R^{r_i}.
\]
The saturated inverse image \(U_{i,s}\) under \(\lambda_i\) has an integral complement \(C_{i,s}\). Thus
\[
 O^{r_i}=U_{i,s}\oplus C_{i,s},\qquad U_{i,s}=\lambda_i^{-1}(U^K_{i,s})\cap O^{r_i}.
\]
Multiplying
the rounded conditions by \(\rho_s(\Lambda)\) shifts the breaks
by \(H_s^{\mathrm{arc}}\). In bases of the saturated submodule
and its complement, Weierstrass division gives exactly the
three divisibility tests in \eqref{eq:wach-level-condition}.

For a constant plus this lattice, subtract the series
\(\Lambda d\) for the corresponding crystalline constant
in \(K_O^{r_i}\) and clear its finitely many coefficient
denominators. This gives the test with \(\varpi^k y-\Lambda d\)
in that paragraph. The crystalline map on the transported
bases has the matrix of the Wach map reduced modulo \(X\), because
reduction, crystalline comparison and Frobenius are functorial.
Thus both evaluation maps commute with the specified maps.
For the crystalline map \(F\), compatibility of the up submodules comes from
\[
 F(\mathcal F_{1,s}^{b_s})\subseteq\mathcal F_{2,s}^{b_s}.
\]

\end{proof}

\paragraph{One integral partner at all levels of an arc.}\label{par:wach-partners}
The truncation modulus after the arc is
\[
 D'=\phi^{a-j}\left(X^M\prod_{s<j}p_s^{3g}\right).
\]
Its orders at the evaluated levels are
\[
 B_s^*=\begin{cases}
 (\ell-1)M,&s<a-j,\\
 3g,&a-j\le s<a .
 \end{cases}
\]

\begin{lemma}[One integral representative for all levels of an arc]\label{lem:wach-arc-partners}
For the cut and arc of the construction, assume \((\ell-1)M>h+2g\),
\(3g\le(\ell-1)M\), and the stated buffer of thin letters.
The comparison uses the filtration
\[
 F_{i,s}=\{z:\operatorname{ord}(z)\ge B_s^*-H_s^{\rm arc}\}.
\] After the scalar loss \(\ell^\sigma\)
of the cut congruences, a vector in either the Wach or the coordinate
description has one integral partner in the other description
agreeing modulo \(F_{i,s}\) at every \(s<a\).
\end{lemma}

\begin{proof}
\proofstep{Orders of the truncation factors.}
The polynomial \(\phi^{a-j}(X)\) has a simple \(X\)-factor and
\(p_s\)-orders \(\ell-1\) for \(s<a-j\). Define \(F_{i,s}\)
to consist of tuples of order at least \(B_s^*-H_s^{\mathrm{arc}}\).
At each of the last \(j\) backward offsets the letter is thin
by the buffer after the starting cut. At thin offsets the order
is at least \(g\), and at up offsets it is at least \(g+1\).

At down offsets it is at least \(b_s+S_0+1\).
Here we use
\[
 (\ell-1)M>h+2g
\]
and \(b_s+S_0\le h\).
All these degrees are in the full-coefficient part of
both filtrations. Thus \(F_{i,s}\) is contained in
\(J_{i,s}\cap B_{i,s}\). The same order is used for the two
representations and the crystalline map preserves it.
If a coordinate evaluation is in \(F_{i,s}\), its
numerator is divisible by \(p_s^{B_s^*}\) by the scalar
remainder test, in particular by \(p_s^{H_s^{\mathrm{arc}}}\).

\proofstep{A basis fixed modulo the cut.}
There is a basis change \(A_i(X)\) modulo \(X^M\), with
constant term the identity, making \(e_i A_i\)
\(\gamma\)-invariant modulo \(X^M\) over \(K_O\).
Construct it degree by degree, dividing at degree \(r\) by
\(\omega_X^r-1\), just as in the cut correction.
Write
\[
 A_i(X)=1+\sum_{r=1}^{M-1}A_{i,r}X^r.
\]

If \(C_{i,r}\) is the degree-\(r\) coefficient of the remaining
invariance defect of the previous truncation in the basis \(e_i\),
the next correction solves
\[
  (\omega_X^r-1)A_{i,r}=-C_{i,r},\qquad
  A_{i,r}\in\ell^{-a_r}\operatorname{Mat}(O),\qquad
  a_r=\sum_{k=1}^r(1+v_\ell(k)).
\]
For the inverse, every term in degree \(r\) of the geometric series
has the form \(\pm A_{i,r_1}\cdots A_{i,r_t}\), where
\(r_1+\cdots+r_t=r\). For a product of coefficients with degrees \(r_1+\cdots+r_t=r\), the extra valuation is controlled by
\[
 \sum_j v_\ell(r_j!)\le v_\ell(r!),\qquad \frac{r!}{\prod_j r_j!}\in\mathbb Z.
\]

In subsequent bounds write \(A_r=a_r\). Use polynomial representatives of the two
truncated matrices. At the root one the ordinary parameter
satisfies \(X\sim t^{\ell-1}\). In crystalline coordinates the
normalized basis at one is constant modulo \(t^{(\ell-1)M}\):
if \(u_k\) is a coefficient in an intermediate degree, cyclotomic invariance gives
\[
 (\omega^k-1)u_k=0,\qquad 0<k<(\ell-1)M ,
\]
and \(\omega^k\ne1\) forces \(u_k=0\). Every evaluated root reaches one after
\(a\) Frobenius steps, which transport this constant to
\(\lambda_i\) and preserve the orders of the nonconstant terms.

\proofstep{The two integral representatives.}
Applying \(\phi^{a-j}\) to the scalar cut congruences shows that
\(\ell^\sigma\phi^a(A_i)\) and
\(\ell^\sigma\phi^a(A_i^{-1})\) are congruent to integral
matrices modulo \(D'\) over \(\mathcal R[1/\ell]\).
Here each coefficient uses the bound with \(A_r\) at its own
degree. Also \(\phi^a(X)^M\) is divisible by \(D'\)
by the upper bound
\[
 3g\le(\ell-1)M.
\]
Consequently
the product of the two transported matrices is the identity
modulo \(D'\) before multiplying by the scalar.
Write an arc vector as \(\Lambda^{-1}(\Phi^a e_i)v\).
Its coordinate partner is obtained by applying an integral
representative of \(\ell^\sigma\phi^a(A_i^{-1})\) to \(v\).

Conversely, the arc partner of a coordinate vector \(y\)
uses \(v\) obtained by applying an integral representative
of \(\ell^\sigma\phi^a(A_i)\) to \(y\).
In either direction the transported invariant basis gives
the required agreement after multiplying by \(\Lambda\), modulo
\(t^{B_s^*}\) at every level. Dividing by \(\Lambda\) gives
the congruence modulo \(F_{i,s}\). Each partner uses a single
integral vector for the whole arc.
\end{proof}

\paragraph{The two cyclotomic polynomial actions.}\label{par:wach-polynomial-actions}
Let \(\omega=\chi(\gamma)\) and let \(\tau_s\) be the permutation
of the primitive roots at level \(s\) given by \(\omega\).
On degree \(k\), \(\gamma\) acts by \(\omega^k\tau_s\).
The principal units act through a cyclic group of order
\(\ell^s\), generated by \(\tau_s\), and together with
the tame units act transitively on the roots.

\begin{lemma}[Polynomial corrections at all evaluated levels]\label{lem:wach-level-polynomials}
On the tame-invariant space,
\[
 (\gamma-1)w_s\in J_{i,s}\quad\Longrightarrow\quad w_s\in D_{i,s}+J_{i,s}.
\] For \(E=a-1\) the polynomials
\(P_{E,I}\) constructed from the respective target and source
degree intervals below send \(p_s^{-g}B_{2,s}\) into
\(D_{2,s}+B_{2,s}\) and \(p_s^{-g}J_{1,s}\) into \(D_{1,s}+J_{1,s}\)
at all offsets. The degrees used lie in at most three intervals
of length \(3g\) inside the prescribed range \([-2h,2h]\).
\end{lemma}

\begin{proof}
Suppose \((\gamma-1)w_s\in J_{i,s}\). In degree \(k\ne0\),
let \(\bar w_{s,k}\) be the fixed class in the degree-\(k\)
quotient. The permutation of the roots has order \(\ell^s\), so
\[
   \bar w_{s,k}=\gamma^{\ell^s}\bar w_{s,k}
               =\omega^{k\ell^s}\bar w_{s,k}.
\]
For \(k\ne0\), \(\omega^{k\ell^s}\ne1\), hence
\(\bar w_{s,k}=0\) over the coefficient field. In degree zero the class is both tame and
principal invariant, thus constant across the roots. A lift
of that constant in \(D_{i,c}\) shows \(w_s\in D_{i,s}+J_{i,s}\).
The action fixes diagonal constants and preserves \(J_{i,s}\).

\proofstep{Averaging the degree-zero coefficient.}
Use the polynomials \(P_{E,I}\) of the arc calculation with
\(E=a-1\), so \(E\ge s\) at each level.
A factor \(Z^{\ell^E}-\omega^{k\ell^E}\) kills degree \(k\).
In degree zero the sum \(\sum_{v=0}^{\ell^E-1}Z^v\) averages
over the principal orbit with a scalar multiplicity.
For a tame-invariant degree-zero tuple \(b\) on the roots, its
\(\gamma\)-orbit has period dividing \(\ell^s\). Hence
\[
 \sum_{v=0}^{\ell^E-1}\gamma^v b
   =\ell^{E-s}\sum_{j=0}^{\ell^s-1}\gamma^j b
   \qquad(E\ge s).
\]
The sum is fixed by \(\gamma\) and by the tame action. Their transitivity
makes it diagonal constant, a space preserved by the remaining factors.

\proofstep{The rounded target and the filtered source.}
For the target polynomial, write \(\mathcal B_s\) for the rounded breaks defining \(B_{2,s}\). Take precisely the integer degrees
\[
 I=\bigcup_{e\in\mathcal B_s}\bigl([e-g,e)\cap\mathbb Z\bigr).
\]

A vector in \(p_s^{-g}B_{2,s}\) satisfies in degree \(k\)
the rounded condition for \(k+g\).
When this condition is larger than that for \(k\), one of
the selected intervals contains \(k\). Every such \(k\)
is nonzero, since \(b_s>g\) on positions with larger jumps.
The polynomial kills these degrees, and in the other degrees
it preserves the required subspace. It fixes the space of
constants up to scalar, giving the required map to
\(D_{2,s}+B_{2,s}\).

On the source side use the separate set
\[
 I_{\rm src}=\bigcup_{u}\bigl([-u-g,-u)\cap\mathbb Z\bigr),
\]
where \(u\) runs over the jumps in the prescribed source boxes.
The degree test for \(p_s^{-g}J_{1,s}\) compares
\(\mathcal F_{1,s}^{-k-g}\) with \(\mathcal F_{1,s}^{-k}\).
Every change occurs in a selected interval. The factors kill
all offending nonzero degrees. In degree zero the principal
sum gives a diagonal constant, which is allowed in
\(D_{1,s}+J_{1,s}\). At degrees where the subspace is
unchanged it is preserved. This proves the second action
on the same evaluation model.

For each polynomial the union lies in three integer intervals
of length at most \(3g\). After enlarging \(G_0\) to include \(\eta_0\), the central interval and its length satisfy
\[
 I_{\rm cent}\subseteq[-G_0-g,G_0-1],\qquad |I_{\rm cent}|\le2G_0+g\le3g.
\]
The intervals around a
thick up or down break lie within the corresponding one-sided
translated window of length at most \(g+S_0\).
The arc contains at most two positions with larger jumps. All degrees lie
in \([-2h,2h]\). These are precisely the window and order
bounds used for the scalar valuation at one. The powers of
\(\ell\) in the partners and polynomial values use only the
maximum length of an arc and the prescribed boxes.
\end{proof}
